\documentclass[11pt]{article}
\usepackage{../../../tex/lecturestyle}

\title{Prüfer Code and Cayley's Formula}
\author{}
\date{}

\begin{document}
\thispagestyle{fancy}

\lecturetitleband{Combinatorics}{Prüfer Code and Cayley's Formula}{6}

\section{Overview}
How many trees can be drawn on $n$ labeled vertices? Writing down a few
cases ($1, 3, 16, 125, \dots$) suggests a pattern, but a direct proof is
awkward, because a tree is a global object: its edges are not independent
of one another. The \emph{Prüfer code} sidesteps the difficulty. It
turns a labeled tree into a plain sequence of $n-2$ labels, with no
constraints at all between the entries, and it does so reversibly.
Counting trees therefore becomes counting sequences.

This lecture defines the code, proves that it is a bijection, and shows a
linear-time implementation. The remaining sections turn the bijection into
a counting toolbox: Cayley's formula, trees with prescribed degrees, ways
to join a forest into a tree, rooted forests with prescribed roots, and
spanning trees of complete bipartite graphs.

\section{Labeled Trees and Cayley's Formula}
\begin{definitionbox}[Labeled tree]
A \emph{labeled tree} on $n$ vertices is a tree whose vertex set is
$\{0,1,\dots,n-1\}$. Two labeled trees are the same exactly when they have
the same set of edges; relabeling the vertices of a drawing gives, in
general, a different labeled tree.
\end{definitionbox}

For $n=3$ there are three labeled trees: a path is determined by which
vertex sits in the middle. For $n=4$ there are $16$: $12$ paths (the
$4!/2$ orderings of the vertices along the path) and $4$ stars (the
choice of the center).

\begin{theorembox}[Cayley's formula]
For $n\ge 2$, the number of labeled trees on $n$ vertices is $n^{\,n-2}$.
\end{theorembox}

The proof comes in Section~\ref{sec:bijection}. The point to notice is the
shape of the answer: $n^{n-2}$ is exactly the number of sequences of
length $n-2$ over an alphabet of size $n$.

\section{The Prüfer Code}
\begin{definitionbox}[Prüfer code]
Let $T$ be a labeled tree on $n\ge 2$ vertices. Repeat $n-2$ times:
\begin{enumerate}[leftmargin=1.5em,itemsep=1pt]
  \item let $u$ be the \emph{smallest-labeled leaf} of the current tree;
  \item append the label of $u$'s unique neighbor to the output;
  \item delete $u$ and its edge.
\end{enumerate}
The output $P=(p_0,\dots,p_{n-3})$ is the \emph{Prüfer code} of $T$.
Two vertices (joined by one edge) are left over and are never written
down. For $n=2$ the code is the empty sequence.
\end{definitionbox}
Each step starts with at least three vertices, so a tree with at least two
leaves is available and ``smallest leaf'' is well defined.

\begin{examplebox}
Take the labeled tree on $9$ vertices in Figure~\ref{fig:pr_tree}.
Table~\ref{tab:pr_trace} lists the deletions. The code is
$P=(4,4,7,7,6,3,6)$; the vertices $6$ and $8$ remain at the end.
\end{examplebox}

% FIGURE: pr_tree
\begin{figure}[H]
\centering
\begin{tikzpicture}[x=1.0cm,y=1.0cm]
\node[lvert,lhi] (n0) at (0.0,1.4) {0};
\node[lvert,lhi] (n1) at (0.0,-0.2) {1};
\node[lvert] (n4) at (1.4,0.6) {4};
\node[lvert,lhi] (n2) at (3.0,1.7) {2};
\node[lvert] (n7) at (3.0,0.6) {7};
\node[lvert] (n3) at (4.6,0.6) {3};
\node[lvert,lhi] (n5) at (6.2,1.7) {5};
\node[lvert] (n6) at (6.2,0.6) {6};
\node[lvert,lhi] (n8) at (7.6,0.6) {8};
\draw[thick,lectureink!70] (n0)--(n4) (n1)--(n4) (n4)--(n7) (n2)--(n7)
  (n7)--(n3) (n3)--(n6) (n5)--(n6) (n6)--(n8);
\end{tikzpicture}
\caption{A labeled tree on nine vertices; shaded vertices are the leaves
$\{0,1,2,5,8\}$. Its Prüfer code is $(4,4,7,7,6,3,6)$.}\label{fig:pr_tree}
\end{figure}

\begin{table}[H]
\centering\small\sffamily
\begin{tabular}{@{}cccc@{}}
\toprule
step $i$ & leaves before the step & deleted leaf $u_i$ & recorded $p_i$\\
\midrule
0 & $\{0,1,2,5,8\}$ & 0 & 4\\
1 & $\{1,2,5,8\}$   & 1 & 4\\
2 & $\{2,4,5,8\}$   & 2 & 7\\
3 & $\{4,5,8\}$     & 4 & 7\\
4 & $\{5,7,8\}$     & 5 & 6\\
5 & $\{7,8\}$       & 7 & 3\\
6 & $\{3,8\}$       & 3 & 6\\
\bottomrule
\end{tabular}
\caption{Encoding the tree of Figure~\ref{fig:pr_tree}. Vertex $4$ becomes a
leaf only after both of its other neighbors $0$ and $1$ have been
deleted.}\label{tab:pr_trace}
\end{table}

Look at the table from the point of view of edges. Every step removes the
edge $\{u_i,p_i\}$, and after the last step the edge $\{6,8\}$ remains.
So the pairs $(u_i,p_i)$ for $i=0,\dots,n-3$, together with the final
pair, are exactly the $n-1$ edges of the tree
(Figure~\ref{fig:pr_tape}). If the deleted leaves $u_i$ could be
recovered from the code alone, the tree could be rebuilt. That is what the
next lemma achieves.

% FIGURE: pr_tape
\begin{figure}[H]
\centering
\begin{tikzpicture}[x=0.95cm,y=1.0cm]
\foreach \i/\u/\p in {0/0/4,1/1/4,2/2/7,3/4/7,4/5/6,5/7/3,6/3/6}{
  \node[lcell,lhi] (u\i) at (\i,1.2) {\u};
  \node[lcell,lhi2] (p\i) at (\i,0) {\p};
  \draw[lthin] (u\i.south)--(p\i.north);
}
\node[lcell,lgray] (uf) at (8.0,1.2) {6};
\node[lcell,lgray] (pf) at (8.0,0) {8};
\draw[lthin] (uf.south)--(pf.north);
\node[llabel,anchor=east] at (-0.6,1.2) {deleted $u_i$};
\node[llabel,anchor=east] at (-0.6,0) {recorded $p_i$};
\node[llabel] at (8.0,-0.7) {last edge};
\node[llabel] at (3.0,-0.7) {code $P$};
\end{tikzpicture}
\caption{Each column is one edge of the tree: the deleted leaf above, the
recorded neighbor below. The Prüfer code is the bottom row; the top row
is what the decoder has to reconstruct.}\label{fig:pr_tape}
\end{figure}

\begin{theorembox}[Degree lemma]
Let $T$ be a labeled tree on $n\ge2$ vertices with Prüfer code $P$. For
every vertex $v$,
\[
  \deg_T(v) \;=\; 1 + \#\{\,i : p_i = v\,\}.
\]
In particular the leaves of $T$ are exactly the labels that do \emph{not}
occur in $P$.

\emph{Proof.} Every edge of $T$ is either the final edge or is deleted at
exactly one step, together with the leaf $u_i$ at that step. Fix $v$ and
look at its incident edges. Either $v$ is deleted at some step, and then
exactly one incident edge (its last one) goes with it; or $v$ is one of the
two survivors, and its last edge is the final edge. Every other incident
edge of $v$ was deleted earlier because its other endpoint was a leaf at
that moment, and that step recorded $v$. Conversely each occurrence of $v$
in $P$ is such an edge. Hence $\deg_T(v)-1$ edges contribute one record
each. \qed
\end{theorembox}

The same lemma applies to every intermediate tree: after $i$ steps the
remaining tree has code $(p_i,\dots,p_{n-3})$, so \emph{a surviving vertex
is currently a leaf exactly when it does not occur in the unread suffix of
the code.}

\section{Decoding and the Bijection}\label{sec:bijection}
The lemma tells a decoder everything it needs. At step $i$ the deleted leaf
$u_i$ is the smallest vertex that has not been deleted yet and does not
occur in $(p_i,\dots,p_{n-3})$, and the edge deleted at that step is
$\{u_i,p_i\}$.

\begin{definitionbox}[Decoding]
Given $P=(p_0,\dots,p_{n-3})\in\{0,\dots,n-1\}^{n-2}$, mark all vertices
unused. For $i=0,1,\dots,n-3$: let $u_i$ be the smallest unused vertex not
in $(p_i,\dots,p_{n-3})$, add the edge $\{u_i,p_i\}$, and mark $u_i$ used.
Finally join the two remaining unused vertices.
\end{definitionbox}

\begin{examplebox}
Decode $P=(3,3,0,5,0,3)$ with $n=8$. The vertices $1,2,4,6,7$ are absent
from $P$. Step $0$ picks $u_0=1$, edge $\{1,3\}$. Step $1$ picks $u_1=2$,
edge $\{2,3\}$. Step $2$ picks $u_2=4$, edge $\{4,0\}$. Now the suffix is
$(5,0,3)$, and $6$ is the smallest unused vertex missing from it:
$u_3=6$, edge $\{6,5\}$. Then $5$ itself no longer occurs in the suffix
$(0,3)$, so $u_4=5$, edge $\{5,0\}$; then $u_5=0$, edge $\{0,3\}$. The
unused vertices $3$ and $7$ are joined last. The result is a tree, and
encoding it gives back $(3,3,0,5,0,3)$.
\end{examplebox}

To make the proof clean, work with an arbitrary finite totally ordered
label set $S$ instead of $\{0,\dots,n-1\}$; the induction below removes a
vertex and therefore leaves the label set.

\begin{theorembox}[Prüfer bijection]
Let $S$ be a totally ordered set with $|S|=n\ge2$. The Prüfer code is a
bijection between trees on the vertex set $S$ and sequences in $S^{\,n-2}$.

\emph{Proof.} Induct on $n$. For $n=2$ there is one tree and one (empty)
sequence. Let $n\ge3$ and fix $P=(p_0,\dots,p_{n-3})\in S^{n-2}$. Fewer
than $n$ distinct values occur in $P$, so
$u=\min\bigl(S\setminus\{p_0,\dots,p_{n-3}\}\bigr)$ exists. Put
$P'=(p_1,\dots,p_{n-3})$; since $u$ does not occur in $P$, we have
$P'\in(S\setminus\{u\})^{n-3}$.

\emph{At most one tree has code $P$.} Suppose $T$ has code $P$. By the
degree lemma the leaves of $T$ are exactly $S\setminus\{p_0,\dots,p_{n-3}\}$,
so the first deleted leaf is $u$, with neighbor $p_0$, and $T-u$ has code
$P'$. By the induction hypothesis $T-u$ is the unique tree on
$S\setminus\{u\}$ with code $P'$; hence $T$ is that tree plus the edge
$\{u,p_0\}$.

\emph{At least one tree has code $P$.} Let $T'$ be the tree on
$S\setminus\{u\}$ with code $P'$ and let $T=T'+\{u,p_0\}$ (note
$p_0\in S\setminus\{u\}$). The leaves of $T$ are $u$ together with the
leaves of $T'$ other than $p_0$. Such a leaf $w\ne p_0$ of $T'$ does not
occur in $P'$ (degree lemma), and $w\ne p_0$, so $w$ does not occur in $P$
at all; by the choice of $u$ we get $w>u$. Thus $u$ is the smallest leaf of
$T$, its neighbor is $p_0$, and deleting it leaves $T'$, whose code is
$P'$. So the code of $T$ is $(p_0,P')=P$. \qed
\end{theorembox}

The proof is the decoding algorithm in disguise: the vertex $u$ chosen in
the induction step is the first $u_0$ of the decoding procedure.

\begin{theorembox}[Cayley's formula]
The number of labeled trees on $n\ge2$ vertices is $n^{\,n-2}$.

\emph{Proof.} By the bijection, it equals the number of sequences in
$\{0,\dots,n-1\}^{n-2}$. \qed
\end{theorembox}

\section{Fast Implementations}
\subsection{With a priority queue: \texorpdfstring{$O(n\log n)$}{O(n log n)}}
\emph{Encoding.} Keep the current degrees and a min-heap of the current
leaves. Pop the smallest leaf $u$; its live neighbor $v$ (the one with
positive degree) is the next code entry; decrease $\deg(v)$, and if it
reaches $1$ push $v$. \emph{Decoding.} Count how often each vertex occurs
in $P$ and keep a min-heap of the vertices with zero remaining occurrences.
For each $p_i$, pop the smallest $u$, output $\{u,p_i\}$, decrease the
count of $p_i$, and push $p_i$ when it reaches zero. Finish by popping the
last two vertices as the final edge.

\subsection{With a moving pointer: \texorpdfstring{$O(n)$}{O(n)}}
The heap is only needed to find ``the smallest current leaf''. Observe how
the leaf set changes in one step: the deleted leaf $u$ disappears, and the
only vertex that can newly become a leaf is its neighbor $p$. Keep a
pointer \texttt{ptr} that scans labels upward and a variable \texttt{leaf}
holding the smallest current leaf. After deleting \texttt{leaf} and
recording $p$:
\begin{itemize}[leftmargin=1.5em,itemsep=1pt]
  \item if $p$ just became a leaf and $p<\texttt{ptr}$, then $p$ is smaller
        than every other leaf (all of those are above \texttt{ptr}), so the
        next leaf to delete is $p$ itself;
  \item otherwise advance \texttt{ptr} to the next label whose degree is
        $1$; this vertex is the new smallest leaf. If $p$ became a leaf
        with $p>\texttt{ptr}$, the scan will reach it later, when its
        degree is already $1$.
\end{itemize}
\texttt{ptr} only moves up, so the total work is $O(n)$. For encoding,
the tree must first be oriented (rooted at $n-1$ is convenient) so that
each leaf knows its remaining neighbor, which is its parent; the root is
never deleted before the last two vertices, because it is the largest
label.

\section{Reference Implementations}
The listings contain all four routines: heap-based and linear-time
encoding and decoding. Both files also cross-check them against each other
and against brute force, as described in
Section~\ref{sec:verify}.

\codenote
\subsection{Rust}
\corelisting{Rust}{example.rs}

\subsection{Python}
\corelisting{Python}{example.py}

\section{Counting with Prüfer Codes}
The degree lemma says that a vertex of degree $d$ appears exactly $d-1$
times in the code. Constraints on degrees therefore become constraints on
the multiset of symbols in a sequence, which is a classical counting
problem.

\subsection{Trees with prescribed degrees}
\begin{theorembox}[Degree sequences]
Let $n\ge2$. A sequence $d=(d_0,\dots,d_{n-1})$ is the degree sequence of
some labeled tree if and only if $d_v\ge1$ for all $v$ and
$\sum_v d_v=2(n-1)$. In that case the number of labeled trees with
$\deg(v)=d_v$ for all $v$ is
\[
  \binom{n-2}{d_0-1,\,d_1-1,\,\dots,\,d_{n-1}-1}
  \;=\;\frac{(n-2)!}{\prod_{v}(d_v-1)!}.
\]

\emph{Proof.} By the bijection and the degree lemma, trees with degree
sequence $d$ correspond to sequences of length $n-2$ in which each label
$v$ occurs exactly $c_v=d_v-1$ times. Such a sequence exists if and only
if $c_v\ge0$ and $\sum c_v=n-2$, which is the stated condition (the sum
condition is also forced by the handshake lemma). The number of such
sequences is the multinomial coefficient. \qed
\end{theorembox}

For instance, on $n=8$ vertices with $d=(4,1,2,1,1,3,1,1)$ the count is
\[
  \frac{6!}{3!\,0!\,1!\,0!\,0!\,2!\,0!\,0!}=60 .
\]

Summing over all degree sequences packages the whole statement into one
identity.
\begin{theorembox}[Degree generating function]
With one indeterminate $x_v$ per vertex,
\[
  \sum_{T}\ \prod_{v=0}^{n-1}x_v^{\deg_T(v)}
  \;=\; x_0x_1\cdots x_{n-1}\,(x_0+x_1+\cdots+x_{n-1})^{\,n-2},
\]
the sum running over all labeled trees on $n$ vertices.

\emph{Proof.} Map each tree to its code $P$; the monomial is
$\prod_v x_v\cdot\prod_{i}x_{p_i}$, and summing $\prod_i x_{p_i}$ over all
free choices of the $n-2$ entries gives $(\sum_v x_v)^{n-2}$. \qed
\end{theorembox}

Setting all $x_v=1$ recovers Cayley's formula. Setting $x_v=t$ for one
vertex and $1$ for the others and comparing coefficients gives the
number of trees in which a given vertex has degree $d$:
\[
  \binom{n-2}{d-1}\,(n-1)^{\,n-1-d}.
\]
(Choose which $d-1$ of the $n-2$ positions hold that vertex; the other
positions take any of the remaining $n-1$ labels.) For $n=4$ and
$d=1,2,3$ this gives $9,6,1$, which sums to $16$.

\subsection{Joining a forest into a tree}
\begin{theorembox}[Forest completion]
Let $F$ be a forest on $n$ vertices with $k\ge 1$ connected components of
sizes $s_1,\dots,s_k$. The number of sets of $k-1$ additional edges that
turn $F$ into a tree is
\[
  n^{\,k-2}\prod_{j=1}^{k}s_j .
\]
(For $k=1$ the value is $1$, the empty set.)

\emph{Proof.} Contract every component to a single vertex $j$ of weight
$s_j$. A valid set of new edges becomes a labeled tree $T$ on
$\{1,\dots,k\}$ (no loops or multi-edges appear, otherwise the result
would contain a cycle). Conversely, given $T$, each of its edges $\{i,j\}$
can be realized by any of $s_is_j$ vertex pairs, independently of the
other edges. So the answer is
\[
  \sum_{T}\prod_{\{i,j\}\in T}s_is_j
  \;=\;\sum_T\prod_{j}s_j^{\deg_T(j)}
  \;=\;\Bigl(\prod_j s_j\Bigr)\Bigl(\sum_j s_j\Bigr)^{k-2},
\]
by the degree generating function with $x_j=s_j$, and $\sum_js_j=n$. \qed
\end{theorembox}

For example, take $7$ vertices split into components of sizes $1$, $2$ and
$4$. The number of ways to connect them with two new edges is
$7^{1}\cdot(1\cdot2\cdot4)=56$.

\subsection{Rooted forests with prescribed roots}
A \emph{rooted forest} is a forest in which one vertex of each component
is designated as its root.

\begin{theorembox}[Prescribed root set]
For $1\le k\le n$, the number of rooted forests on $\{0,\dots,n-1\}$ whose
set of roots is exactly $R=\{0,\dots,k-1\}$ equals $k\,n^{\,n-k-1}$.

\emph{Proof.} Add a new vertex $n$ and join it to every root. This is a
bijection between such forests and labeled trees on $\{0,\dots,n\}$ in
which the neighbor set of $n$ is exactly $R$: deleting $n$ from such a tree
leaves a forest in which each component contains exactly one vertex of
$R$ (two would close a cycle through $n$, none would disconnect the
tree). By symmetry among labels, every $k$-subset of $\{0,\dots,n-1\}$ is
the neighbor set of $n$ equally often. Trees on $n+1$ vertices where $n$
has degree $k$ number $\binom{n-1}{k-1}n^{\,n-k}$ (previous subsection),
so the answer is
\[
  \binom{n-1}{k-1}n^{\,n-k}\Big/\binom{n}{k}\;=\;k\,n^{\,n-k-1}. \qed
\]
\end{theorembox}

For $k=1$ this is Cayley's formula again, since a rooted tree with
prescribed root $0$ is just a tree.

\subsection{Spanning trees of complete bipartite graphs}
\begin{theorembox}[$K_{a,b}$]
The complete bipartite graph $K_{a,b}$ has $a^{\,b-1}b^{\,a-1}$ spanning
trees.

\emph{Proof sketch.} Let the parts $A$ and $B$ have sizes $a$ and $b$, and
run the Prüfer procedure on a spanning tree (labels of the two parts may
be interleaved arbitrarily), stopping when one vertex of each part is left.
A deleted leaf in $A$ has its neighbor in $B$, and vice versa. Store each
recorded neighbor in the sequence belonging to \emph{its own part}: this
yields $P_A\in A^{b-1}$ (one entry per deleted vertex of $B$) and
$P_B\in B^{a-1}$.

The degree lemma still holds with this split: a vertex of $A$ has degree
one more than its number of occurrences in $P_A$, and likewise for $B$.
Decoding mimics the ordinary one: the next leaf $u$ is the smallest unused
vertex absent from the unread part of the sequence of its own part, and
its neighbor is the next unread entry of the \emph{other} sequence. The
procedure never gets stuck: when $a'$ vertices of $A$ and $b'$ vertices of
$B$ are unused, the unread parts of $P_A$ and $P_B$ have lengths $b'-1$ and
$a'-1$, so a vertex of the larger side (or either, when $a'=b'$) is always
a leaf, and the other sequence is non-empty until $a'=b'=1$. The induction
used for the bijection theorem goes through unchanged, so spanning trees
correspond to pairs $(P_A,P_B)$, and there are $a^{b-1}\cdot b^{a-1}$ of
them. \qed
\end{theorembox}

As a sanity check, $K_{2,3}$ has $2^2\cdot3^1=12$ spanning trees and
$K_{3,3}$ has $3^2\cdot3^2=81$. The harness in \texttt{example.py} also
implements the two-sequence code and checks that it is a bijection.

\section{Verification}\label{sec:verify}
Both reference implementations check the following by brute force.
\begin{itemize}[leftmargin=1.5em,itemsep=1pt]
  \item The tree of Figure~\ref{fig:pr_tree} encodes to
        $(4,4,7,7,6,3,6)$, and $(3,3,0,5,0,3)$ decodes to the edges
        $\{0,3\},\{0,4\},\{0,5\},\{1,3\},\{2,3\},\{3,7\},\{5,6\}$.
  \item Enumerating all edge subsets confirms that there are
        $1,3,16,125,1296$ labeled trees for $n=2,\dots,6$ (both files) and
        $16807$ for $n=7$ (Python), and that encoding is injective on
        them (Python).
  \item For every sequence in $\{0,\dots,n-1\}^{n-2}$ with $n\le7$,
        decoding yields a tree whose code is that sequence.
  \item On hundreds of random trees with up to a few hundred vertices, the heap and
        linear encoders agree, decoding returns the original tree, and
        $\deg(v)=1+\#\{i:p_i=v\}$ holds (the random trees are generated
        without using the code, so the test is not circular).
  \item Python only: the multinomial degree-sequence count for all $n\le7$,
        the single-vertex degree distribution, the forest-completion count
        on sampled forests (including the sizes $1,2,4$ above, giving
        $56$), the prescribed-roots count for all $k\le n\le7$, and the
        $K_{a,b}$ counts together with the two-sequence code.
\end{itemize}

\section{Using the Technique}
\begin{itemize}[leftmargin=1.5em,itemsep=2pt]
  \item \textbf{Count or sum over trees with degree constraints.} Pass to
        sequences: a vertex of degree $d$ is a symbol used $d-1$ times. A
        weight $\prod_v w_v^{\deg(v)}$ becomes the generating function
        above.
  \item \textbf{Connect a given forest.} Use the forest-completion formula;
        only the component sizes matter. Remember that $k=1$ gives $1$.
  \item \textbf{Generate or reconstruct a tree.} A random uniform labeled
        tree is a random sequence plus the $O(n)$ decoder; conversely a
        tree can be stored as $n-2$ integers.
  \item \textbf{Modular arithmetic.} Powers like $n^{k-2}$ with $k=1$ need
        the modular inverse of $n$ (or a special case); multinomials need
        factorials and inverse factorials up to $n$.
\end{itemize}
Problems on the same theme, roughly from easier to harder:
\begin{itemize}[leftmargin=1.5em,itemsep=1pt]
  \item Luogu P6086 --- implement encoding and decoding.
  \item Codeforces 156D \emph{Clues} --- forest completion.
  \item Codeforces 1109D \emph{Sasha and Interesting Fact from Graph
        Theory} --- trees with a prescribed path between two vertices.
  \item Codeforces 917D \emph{Stranger Trees} --- trees containing a
        given forest, via inclusion--exclusion.
  \item AtCoder ABC303 Ex \emph{Constrained Tree Degree} and AtCoder
        ARC106 F \emph{Figures} --- degree-sequence counting.
\end{itemize}

\section{Further Reading}
\begin{itemize}[leftmargin=1.5em]
  \item AtCoder Algorithm Lectures index (source of this curriculum):
        \url{https://info.atcoder.jp/entry/algorithm_lectures/index}
  \item H.\ Prüfer, ``Neuer Beweis eines Satzes über Permutationen'',
        \emph{Archiv der Mathematik und Physik} 27 (1918), pp.\ 142--144.
  \item A.\ Cayley, ``A theorem on trees'', \emph{Quarterly Journal of
        Pure and Applied Mathematics} 23 (1889), pp.\ 376--378.
  \item M.\ Aigner and G.\ M.\ Ziegler, \emph{Proofs from THE BOOK},
        chapter ``Cayley's formula for the number of trees'' (four
        different proofs).
  \item cp-algorithms, \emph{Prüfer code}:
        \url{https://cp-algorithms.com/graph/pruefer_code.html}
\end{itemize}

\end{document}
